\documentclass[../../main.tex]{subfiles}
\begin{document}

\begin{center}
  {\fontsize{48pt}{58pt}\selectfont \textbf{1975}\normalsize\textbf{年度}}
\end{center}

\begin{table}[htb]
  \centering
  \caption{出題テーマと難易度}
  \begin{tabular}{|c|p{9cm}|c|}
    \hline
    \textbf{問題番号} & \textbf{テーマ} & \textbf{難易度} \\
    \hline
    第1問 & 総当たり戦の順位と勝敗 & \\
    \hline
    第2問 & 平方根の評価と定積分の不等式 & \\
    \hline
    第3問 & 三角関数の値が等比数列になる条件 & \\
    \hline
    第4問 & 三角形の重心と円周上の最遠点 & \\
    \hline
    第5問 & 逆数の和を含む式の最小化 & \\
    \hline
    第6問 & 三角形と円の共通部分の面積 & \\
    \hline
  \end{tabular}
\end{table}

\begin{table}[htb]
  \centering
  \caption{解答}
  \renewcommand{\arraystretch}{1.4}
  \begin{tabular}{|c|c|p{9cm}|}
    \hline
    \textbf{問題番号} & \textbf{小問} & \textbf{答え} \\
    \hline
    第1問 & -- & 上位のチームは下位の全チームに勝っている（証明問題） \\
    \hline
    第2問 & (1) & $\sqrt{x}\le\dfrac{x+1}{2}$（証明問題） \\
    \cline{2-3}
    & (2) & 指定された積分は$1$未満（証明問題） \\
    \hline
    第3問 & -- & $\alpha=\beta-k\pi,\gamma=\beta+k\pi$（$k\in\mathbb{Z}$） \\
    \hline
    第4問 & -- & 三角形の重心を$G$，円の中心を$O$とすると，$G\ne O$なら直線$OG$上で$G$から最も遠い円周上の点．$G=O$なら任意の点 \\
    \hline
    第5問 & (1) & $(a_1,a_2,a_3)=(2,3,7)$ \\
    \cline{2-3}
    & (2) & $n=3$，$(a_1,a_2,a_3)=(2,3,7)$．最小値は$\dfrac{1}{42}$ \\
    \hline
    第6問 & (i) & $1\le x<\sqrt{2}$ \\
    \cline{2-3}
    & (ii) & $S=\pi-2F(\dfrac{x}{\sqrt{2}})-F(\sqrt{2}-x)$ \\
    \cline{2-3}
    & (iii) & $x=\dfrac{3\sqrt{2}}{4}$で最大 \\
    \hline
  \end{tabular}
\end{table}

\paragraph{試験形態}
6問．試験時間は要確認．

\end{document}
